\documentclass[11pt,a4paper]{article}
\usepackage[utf8]{inputenc}
\usepackage[T1]{fontenc}
\usepackage[spanish]{babel}
\usepackage{hyperref}
\usepackage{geometry}
\geometry{margin=2.5cm}
\usepackage{parskip}

\title{324 el otro mueve la vista, no el tiempo}
\author{Alejandro Valenzuela Rivera\\\small Universidad Autónoma de Nuevo León}
\date{\today}
\begin{document}
\maketitle

\begin{quote}
\textbf{Sesgo declarado.} Este trabajo se sostiene desde el sesgo declarado de su autor, y no finge neutralidad. Lo que afirma se puede tumbar, y ese es el punto: el falsificador está escrito, no escondido.
\end{quote}

\section*{Resumen}
Doce años de ver sin ser contradicho no moderan una postura: la arman. Lo que la mueve es otro que la contradiga.

\section{Por qué se sostiene}
La regla se sostiene mientras alguien la pueda contestar con lo que Malcolm X (1964, Hajj) muestra.

\section{De dónde viene}
Tesis/referencias/index.md, y Malcolm X (1964, Hajj).

\section{El caso}
Malcolm X: doce años de ver sin ser contradicho lo armaron; ver a otros lo desarmó.

\section{Método}
El tiempo no corrige: solo acumula. Lo que corrige es una mirada que no es la propia.

\section{Límites, y qué lo tumba}
Un caso donde el tiempo solo corrija una postura.

\section{Lo que queda abierto}
\textquestiondown{}Qué la movió: el tiempo, o alguien que la contradijo?

\section*{Cómo se leyó esta tesis}
Esto se escribió en el repositorio man-made-empathy, del que este candado es una pieza: candados/06 del caso/324 el otro mueve la vista, no el tiempo.md. La máquina que lo compuso lo hizo desde este mismo archivo, y no agregó nada que no estuviera escrito. Si algo de esto es falso, se dice en el mismo sitio, y se corrige con fecha.

\end{document}
